\documentclass[../../../include/open-logic-section]{subfiles}

\begin{document}

\olfileid[id]{sth}{story}{rus}
\olsection{Paradoks Russell (Lagi)}

Dalam \olref[sfr][][]{part}, kita bekerja dengan teori himpunan naif. Namun,
menurut suatu konsepsi yang \emph{sangat} naif, himpunan hanyalah ekstensi
predikat. Gagasan naif ini mengharuskan prinsip berikut:

\begin{defish}
  \emph{Komprehensi Naif.} $\Setabs{x}{\phi(x)}$ ada untuk setiap formula $\phi$.
\end{defish}

Walaupun menggoda, prinsip ini dapat dibuktikan tidak konsisten. Kita telah
melihatnya dalam \olref[sfr][set][rus]{sec}, tetapi hasil ini begitu penting
dan begitu langsung sehingga layak diulangi. Apa adanya.

\begin{thm}[Paradoks Russell]
Tidak ada himpunan $R = \Setabs{x}{x \notin x}$
\end{thm}

\begin{proof}
Jika $R = \Setabs{x}{x \notin x}$ ada, maka
$R \in R$ jika dan hanya jika $R \notin R$, yang merupakan kontradiksi.
\end{proof}

Russell menemukan hasil ini pada Juni 1901. (Namun, ia tidak menyajikan
paradoks itu tepat dalam bentuk yang baru saja kita gunakan, sebab ia sedang
mempertimbangkan teori himpunan Frege sebagaimana diuraikan dalam
\emph{Grundgesetze}. Kita akan kembali ke hal ini dalam \olref[blv]{sec}.)
Russell menulis surat kepada Frege pada 16 Juni 1902 untuk menjelaskan
ketidakkonsistenan dalam sistem Frege. Untuk korespondensi tersebut dan
sedikit latar belakang, lihat \citet[pp.~124--8]{Heijenoort1967}.

Perlu ditekankan bahwa bukti dua baris ini merupakan hasil \emph{logika
murni}. Memang, dalam notasi $\Setabs{x}{x \notin x}$ kita secara tersirat
menggunakan sebuah aksioma (nonlogis?), yaitu Ekstensionalitas; sebab
$\Setabs{x}{\phi(x)}$ dimaksudkan sebagai \emph{satu-satunya} himpunan
(berdasarkan Ekstensionalitas) dari objek-objek yang memenuhi $\phi$, jika
himpunan seperti itu ada. Namun, bahkan kesan penggunaan Ekstensionalitas
dapat kita hindari dengan menyatakan hasilnya sebagai berikut:
\emph{tidak ada himpunan yang anggotanya tepat semua himpunan yang bukan
anggota dirinya sendiri}. Dan hal ini tidak banyak berkaitan dengan himpunan.
Sebagaimana diamati Russell sendiri, penalaran yang persis serupa akan
membawa Anda pada kesimpulan: \emph{tidak ada pria yang mencukur tepat semua
pria yang tidak mencukur dirinya sendiri}. Atau: \emph{tidak ada anjing pug
yang mengendus tepat semua anjing pug yang tidak mengendus dirinya sendiri}.
Dan seterusnya. Secara skematis, bentuk hasilnya hanyalah:
\[
\lnot \exists x \forall z(Rzx \liff \lnot R zz).
\]
Dan itu hanyalah sebuah (skema) teorema logika orde pertama. Akibatnya, kita
tidak dapat menghindari Paradoks Russell hanya dengan mengutak-atik teori
himpunan kita; paradoks itu muncul bahkan sebelum kita \emph{sampai} pada
teori himpunan. Jika kita hendak menggunakan logika orde pertama (klasik),
kita harus \emph{menerima} bahwa tidak ada himpunan
$R = \Setabs{x}{x\notin x}$.

Kesimpulannya ialah sebagai berikut. Jika Anda ingin menerima Komprehensi
Naif sembari \emph{menghindari} ketidakkonsistenan, Anda tidak dapat sekadar
mengutak-atik \emph{teori himpunannya}. Sebaliknya, Anda harus merombak
\emph{logikanya}.

Tentu saja, teori-teori himpunan dengan logika nonklasik pernah diajukan.
Namun, teori-teori itu---paling tidak---tidak standar. Pendekatan standar
terhadap Paradoks Russell ialah memperlakukannya sebagai bukti langsung atas
ketiadaan, lalu mencoba mempelajari cara hidup dengannya. Itulah pendekatan
yang akan kita ikuti.

\end{document}
